% Generated by scripts/assemble.py; corrections belong in source/corrections.json.
\examhead{169}{电子科技大学 2026 年数学分析试卷}{160}
% p161-b23; source PDF page 161
\par\Needspace{5\baselineskip}
\blockanchor{p161-b23}
\majorquestion{一}填空题. 每题 5 分, 共 30 分.\par
% p161-b24; source PDF page 161
\blockanchor{p161-b24}
\question{1}设 $a > 0$, 求 $\lim_{n \to \infty} n(\sqrt[n]{a} - \sqrt[2n]{a}) =$ \_\_\_\_.\par
% p161-b25; source PDF page 161
\blockanchor{p161-b25}
\question{2}求  $\lim_{x\to+\infty}\frac{1}{x}\int_{\frac{1}{x}}^{1}\frac{\cos2t}{t^{2}}dt=$  \_\_\_\_.\par
% p161-b26; source PDF page 161
\blockanchor{p161-b26}
\question{3}求 $\int \sqrt{x} e^{\sqrt{x}} \mathrm{d}x =$ \_\_\_\_.\par
% p162-b00; source PDF page 162
\blockanchor{p162-b00}
\question{4}令 $u = \ln (x + y)$, 则 $\mathrm{d}u|_{(1,2)} =$ \_\_\_\_.\par
% p162-b01; source PDF page 162
\blockanchor{p162-b01}
\question{5}求  $\sum_{n=1}^{\infty}\frac{n!}{n^{n}}x^{n}$  的收敛半径 \_\_\_\_.\par
% p162-b02; source PDF page 162
\blockanchor{p162-b02}
\question{6}求  $\int_{0}^{+\infty}\frac{\cos x - \cos 2x}{x^{2}}dx=$  \_\_\_\_.\par
% p162-b03; source PDF page 162
\par\Needspace{5\baselineskip}
\blockanchor{p162-b03}
\majorquestion{二}计算题. 每题 12 分, 共 48 分.\par
% p162-b04; source PDF page 162
\blockanchor{p162-b04}
\question{7}已知 $f(x)$ 在 $[0, +\infty)$ 上连续，且 $\lim_{x \to +\infty} f(x) = A$，求 $\lim_{n \to \infty} \int_0^1 f(nx) \, \mathrm{d}x$.\par
% p162-b05; source PDF page 162
\blockanchor{p162-b05}
\question{8}设 $f(t) = \iint_{0 \leq x \leq t, 0 \leq y \leq t} e^{-\frac{t\pi}{y^2}} \, \mathrm{d}x \, \mathrm{d}y, t > 0$, 求 $f'(t)$.\par
% p162-b06; source PDF page 162
\par\Needspace{6\baselineskip}
\blockanchor{p162-b06}
\question{9}已知曲面\par
\nopagebreak[4]
% p162-b07; source PDF page 162
\blockanchor{p162-b07}
\displayformula{\left\{ \begin{array}{l} x = \frac {1}{2} (e ^ {u} + e ^ {- u}) \cos v; \\ y = \frac {1}{2} (e ^ {u} + e ^ {- u}) \sin v; \\ z = \frac {1}{2} (e ^ {u} - e ^ {- u}). \end{array} \right.}
% p162-b08; source PDF page 162
\blockanchor{p162-b08}
求曲面在点 $(u, v) = \left(0, \frac{\pi}{4}\right)$ 处的切平面和法线方程.\par
% p162-b09; source PDF page 162
\par\Needspace{6\baselineskip}
\blockanchor{p162-b09}
\question{10}求曲面积分\par
\nopagebreak[4]
% p162-b10; source PDF page 162
\blockanchor{p162-b10}
\displayformula{\iint_ {\Sigma} \frac {1}{x} \mathrm{d} y \mathrm{d} z + \frac {1}{y} \mathrm{d} z \mathrm{d} x + \frac {1}{z} \mathrm{d} x \mathrm{d} y.}
% p162-b11; source PDF page 162
\blockanchor{p162-b11}
其中  $\Sigma$  为椭球面  $\frac{x^{2}}{a^{2}} + \frac{y^{2}}{b^{2}} + \frac{z^{2}}{c^{2}} = 1$ ，方向取外侧.\par
% p162-b12; source PDF page 162
\par\Needspace{5\baselineskip}
\blockanchor{p162-b12}
\majorquestion{三}证明题. 每题 12 分, 共 36 分.\par
% p162-b13; source PDF page 162
\blockanchor{p162-b13}
\question{11}设 $\mathbb{R}$ 为实数集, $F \subset \mathbb{R}$ 为非空闭集, 定义 $d(x) = \inf_{t \in F} |x - t|$, 证明 $d(x)$ 为连续函数.\par
% p162-b14; source PDF page 162
\blockanchor{p162-b14}
\question{12}设 $f(x) \in R[0,1]$, 且 $f$ 在 $x = 0$ 处连续, 证明: $\lim_{n \to \infty} \int_0^1 f(x^n) \, \mathrm{d}x = f(0)$.\par
% p162-b15; source PDF page 162
\blockanchor{p162-b15}
\question{13}$x \in [-\pi, \pi]$, 判断 $\sum_{n=2}^{\infty} \frac{\cos nx}{\ln n}$ 的敛散性(绝对, 条件, 发散), 并证明其不是某 Riemann 可积函数在 $[- \pi, \pi]$ 上的 Fourier 级数.\par
% p162-b16; source PDF page 162
\par\Needspace{5\baselineskip}
\blockanchor{p162-b16}
\majorquestion{四}综合题. 每题 18 分, 共 36 分.\par
% p162-b17; source PDF page 162
\blockanchor{p162-b17}
\question{14}设 $K$ 是一个非空紧集, $\{f_n(x)\}$ 为 $K$ 上的连续函数列, 且对任意的 $x \in K$, $\{f_n(x)\}$ 关于 $n$ 单调递减趋于 0, 证明: $\{f_n(x)\}$ 在 $K$ 上一致收敛于 $f(x) = 0$.\par
% p162-b18; source PDF page 162
\blockanchor{p162-b18}
\question{15}设 $\Omega$ 为 $\mathbb{R}^2$ 上的光滑简单连通区域, $f(x) \in C^2(\Omega)$, 若 $\frac{\partial^2 f}{\partial x_1^2} + \frac{\partial^2 f}{\partial x_2^2} = 0$, $\forall x = (x_1, x_2) \in \Omega$, 则称 $f$ 为 $\Omega$ 上的调和函数.\par
% p162-b19; source PDF page 162
\blockanchor{p162-b19}
(1) 证明: \displayformula{\iint_{\Omega} \left( \frac{\partial^2 f}{\partial x_1^2} + \frac{\partial^2 f}{\partial x_2^2} \right) \mathrm{d}\sigma = \int_{\partial \Omega} \frac{\partial f}{\partial n} \mathrm{d}s}, 其中 $n$ 是 $\partial \Omega$ 的单位外法向量.\par
% p162-b20; source PDF page 162
\par\Needspace{6\baselineskip}
\blockanchor{p162-b20}
(2) 证明: $f$ 为 $\Omega$ 上的调和函数当且仅当\par
\nopagebreak[4]
% p162-b21; source PDF page 162
\blockanchor{p162-b21}
\displayformula{f (\boldsymbol {x}) = \frac {1}{2 \pi r} \int_ {\partial B (\boldsymbol {x}, r)} f (\boldsymbol {x}) \mathrm{d} s = \frac {1}{\pi r ^ {2}} \iint_ {B (\boldsymbol {x}, r)} f (\boldsymbol {x}) \mathrm{d} \sigma .}
% p162-b22; source PDF page 162
\blockanchor{p162-b22}
其中  $B(x,r)$  是以 x 为圆心以 r 为半径的闭球.\par
